\subsection{Topologies defined by semi-norms}

Let \(p\) be a semi-norm on the vector space \(E\) over \(K\) ; for every \(\alpha > 0\) let \(V(p, \alpha)\) be the subset of those \(x\) of \(E\) for which \(p(x) \leqslant \alpha\) . Clearly, if \(x \in V(p, \alpha)\) and \(\lambda \in K\) is such that \(|\lambda| \leqslant 1\) , then \(\lambda x \in V(p, \alpha)\) , in other words \(V(p, \alpha)\) is balanced. Further, for every \(x_0 \in E\) , there exists a non-zero scalar \(\mu \in K\) such that \(|\mu| \geqslant p(x_0) \alpha^{-1}\) , therefore \(\mu^{-1}x_0 \in V(p, \alpha)\) that is to say \(V(p, \alpha)\) is absorbent. Finally, from \((\mathrm{SN}_{\mathrm{II}})\) , we have \(V(p, \alpha/2) + V(p, \alpha/2) \subset V(p, \alpha)\) , and from \((\mathrm{SN}_{\mathrm{I}})\) that for every non-zero scalar \(\lambda\) in \(K\) we have \(\lambda V(p, \alpha) = V(p, |\lambda| \alpha)\) . We conclude from these remarks, by I, p.~7, prop. 4, that, when \(\alpha\) varies in the set of numbers \(> 0\) (or only in a sequence of strictly positive numbers tending to 0) then the sets \(V(p, \alpha)\) constitute a fundamental system of neighbourhoods of 0 for a topology compatible with the vector space structure of \(E\) ; we say that this topology is defined by the semi-norm \(p\) . A vector space \(E\) with such a topology is called a semi-normed space. Note that if \(W(p, \alpha)\) is the subset of \(x\) of \(E\) such that \(p(x) < \alpha\) , then the \(W(p, \alpha)\) constitute (where \(\alpha > 0\) , or \(\alpha\) )

varies in a strictly positive sequence of numbers tending to zero) a fundamental system of neighbourhoods of 0 for the topology defined by \(p\) .

If \(\Gamma\) is a set of semi-norms on \(\mathbf{E}\) , then the upper bound of the topologies defined by the semi-norms \(p \in \Gamma\) is compatible with the vector space structure (I, p.~10, cor. 4). A fundamental system of neighbourhoods of 0, for this topology, is given by the finite intersections \(\bigcap_{i} \mathrm{V}(p_i, \alpha_i)\) where \(p_i \in \Gamma\) and \(\alpha_i > 0\) . This topology is said to be defined

by the set of semi-norms \(\Gamma\) . It is the coarsest topology on \(\mathbf{E}\) amongst those that are invariant under all translations and for which the semi-norms \(p \in \Gamma\) are continuous.

Let E be a topological vector space over K : a system of semi-norms on E, say \(\Gamma\) , is called a fundamental system of semi-norms if the topology on E is the same as the topology defined by \(\Gamma\) .

Let E be a vector space over K, with the topology defined by a set of semi-norms \(\Gamma\) . For every semi-norm p, we have \(p(x - z) \leqslant p(x - y) + p(y - z)\) , which shows that the function \((x, y) \mapsto p(x - y)\) is a pseudometric on E (GT, IX, §1.1): it follows from the definitions that, when p varies in \(\Gamma\) , the set of these pseudometrics defines the uniform structure of the topological vector space E.

Remarks. --- 1) The topology defined by a finite set of semi-norms \(p_i\) ( \(1 \leqslant i \leqslant n\) ) on E, can be defined by the single semi-norm \(p = \sup_{1 \leqslant i \leqslant n} p_i\) . But a topology defined by

an infinite set of semi-norms cannot, in general, be defined by a single semi-norm (III, p.~37, exerc. 2).

\begin{enumerate}
\def\labelenumi{\arabic{enumi})}
\setcounter{enumi}{1}
\item
  Let \((\mathcal{T}_{\iota})_{\iota\in I}\) be a family of topologies on a vector space E over K, each of which is defined by a family of semi-norms \(\Gamma_{\iota}\) . Then the topology defined by the set of semi-norms \(\Gamma = \bigcup \Gamma_{\iota}\) is the upper bound of the topologies \(T_{\iota}\) .
\item
  If \(\Gamma_0\) is a set of semi-norms directed by the increasing order relation defined between two semi-norms \(p, q\) on E by « there exists \(\lambda > 0\) such that \(p \leqslant \lambda q\) », then a fundamental system of neighbourhoods of 0, for the topology defined by \(\Gamma_0\) , is obtained by taking the sets \(\mathbf{V}(p, \alpha)\) where \(p \in \Gamma_0\) and \(\alpha > 0\) . If \(\Gamma\) is any set of semi-norms on E, then a filtered set of semi-norms, defining the same topology as \(\Gamma\) , is the set \(\Gamma_0\) of upper envelopes of all finite families of semi-norms belonging to \(\Gamma\) .
\item
  Even if \(\mathbf{K} = \mathbf{R}\) , the topology of a topological vector space over \(\mathbf{K}\) cannot always be defined by a set of semi-norms (cf.~II, p.~24, corollary).
\end{enumerate}

Example. --- Let \(\mathcal{C}^{\infty}(\mathbf{R})\) be the vector space over R of real valued functions that are infinitely differentiable in R. For every function and every pair of integers \(n \geqslant 0\) , \(m \geqslant 1\) , put

\[
p _ {n, m} (f) = \sup _ {- m \leqslant t \leqslant m} | f ^ {(n)} (t) |\tag{2}
\]

with \(f^{(0)} = f\) . Obviously the \(p_{n,m}\) are semi-norms on \(\mathcal{C}^{\infty}(\mathbf{R})\) . In order that the functions \(f_{\alpha}\) tend to 0 (following a filter \(\mathfrak{F}\) on the set of indices) in \(\mathcal{C}^{\infty}(\mathbf{R})\) for the topology \(\mathcal{T}\) defined by the semi-norms \(p_{n,m}\) , it is necessary and sufficient that for all integers \(n \geqslant 0\) , the functions \(f_{\alpha}^{(n)}\) tend to 0 (following \(\mathfrak{F}\) ) uniformly on every compact subset of \(\mathbf{R}\) . We say that \(\mathcal{T}\) is the topology of compact convergence for the functions \(f \in \mathcal{C}^{\infty}(\mathbf{R})\) and all their derivatives (cf.~III, p.~9).

PROPOSITION 2. --- On a vector space E, let T be the topology defined by a set of semi-norms Γ.

\begin{enumerate}
\def\labelenumi{(\roman{enumi})}
\item
  The closure of \(\{0\}\) in \(\mathbf{E}\) , for \(\mathcal{T}\) , is the subset of \(x \in \mathbf{E}\) for which \(p(x) = 0\) for every semi-norm \(p \in \Gamma\) .
\item
  If \(\mathcal{T}\) is Hausdorff and \(\Gamma\) is enumerable, then \(\mathcal{T}\) is metrisable.
\end{enumerate}

The proposition follows immediately from the definitions and from GT, IX, § 2.4, cor. 1.

Note that if \(\mathcal{T}\) is metrisable, it may be that \(\mathcal{T}\) cannot be defined by a single norm; this is the case in the example given above (cf.~IV, p.~18, Example 4).

Let E be a vector space over K, with the topology defined by a set of semi-norms \(\Gamma\) . Let \(\hat{E}\) be the Hausdorff completion of E (I, p.~6), and \(\hat{\Gamma}\) be the set of mappings \(\hat{p}\) of \(\hat{E}\) in \(R_{+}\) where p varies in \(\Gamma\) (GT, II, § 3.7, prop. 15). By the principle of extending inequalities, the functions \(\hat{p} \in \hat{\Gamma}\) are semi-norms on \(\hat{E}\) , and the functions \(\hat{p}(x - y)\) form a set of pseudometrics defining the uniform structure of \(\hat{E}\) (GT, IX, § 1.3, prop. 1). We see, therefore, that \(\hat{\Gamma}\) is a fundamental set of semi-norms defining the topology of \(\hat{E}\) .

